\item En un experimento de anillos de Newton, una lente plano-convexa ($n = 1.52$) con un diámetro de $5.00\text{ cm}$ se coloca sobre una placa plana de vidrio. Con luz de $\lambda = 650\text{ nm}$ a incidencia normal, se cuentan exactamente $55$ anillos brillantes, quedando el último justo en el borde de la lente.
\begin{enumerate}
    \item ¿Cuál es el radio de curvatura $R$ de la superficie esférica de la lente?
    \item ¿Cuál es su distancia focal?
\end{enumerate}